\section{A lower bound for the slope across the components}
\label{sec:cov-heights}

The geometric argument compares heights of individual components with the
arithmetic slope of the space of sections on their union. Both the error in
their relative degrees and the error in their ranks enter this comparison.
We give the finite-sum argument used in the source, including the constant
with which it enters the covolume estimate.

Fix a nonempty finite set \(I\) of components. Write \(h_c\) for a component
height, \(s_c\) for its comparison slope, \(\delta_c\) for its geometric
degree, \(d_c\) for its rank and \(a_c\) for its arithmetic degree. Let
\(D\ge0,\ T\ge3D+5,\ A>0,\ B>0,\ \delta_*>0,\ E\ge1\). The
parameters \(u,v,X,R\) are nonnegative, and \(0\le z,\rho<1\).
Assume at every \(c\in I\) that
\begin{align}
 \delta_*\le\delta_c\le\delta_* /(1-\rho),\qquad
 B\delta_c\le d_c,\qquad (1-z)d_c\le B\delta_c, \label{eq:cov-shares}\\
 h_c+u\ge0,\qquad
 (1-\rho)\max(0,s_c+u)\le h_c+u+v, \label{eq:cov-shift}\\
 a_c\ge B\delta_c
 \bigl(Th_c-3(D+1)|h_c|-A\log(2+|h_c|)-R\bigr).\label{eq:cov-hs}
\end{align}
Assume also \(|I|^{-1}\sum s_c\ge-X\). If
\begin{align}
 TX\le C_1 E,\qquad Tv\le C_1 E,\qquad
 u\bigl(6D+10+T(z+2\rho)\bigr)\le C_1E, \nonumber\\
 A\log(\max(A,1))+R\le C_2 E,\qquad
 2\frac{\sum a_c}{\sum d_c}\le\mu,\label{eq:cov-budgets}
\end{align}
then
\begin{equation}
                  \mu\ge-2(3C_1+C_2+4)E. \label{eq:cov-u}
\end{equation}

\paragraph{Proof.}
Set \(m=\max(0,(1-\rho)u-X-v)\) and \(n=|I|\). Averaging
\eqref{eq:cov-shift} gives
\[
 \frac1n\sum(h_c+u)\ge (1-\rho)(u-X)-v
 \ge (1-\rho)u-X-v.
\]
The left side is also nonnegative, so it is at least \(m\).
For \(t\ge0\) one has
\[
 A\log(2+t)\le t+A\log(\max(A,1))+2 .
\]
If \(A\ge1\), apply \(\log y\le y-1\) to \(y=(2+t)/A\). If
\(0<A<1\), use \(\log(2+t)\le1+t\) and \(A(1+t)\le1+t\).
Put \(P=T-3D-4\), \(K=A\log(\max(A,1))+2\), and
\[
 N=Pu+2u(3D+4)+K+R,\qquad W=\sum B\delta_c,\qquad S=\sum d_c.
\]
Since \(h_c\ge-u\), one has \(|h_c|\le h_c+2u\). Substituting both
inequalities in \eqref{eq:cov-hs} gives
\[
 a_c\ge B\delta_c\bigl(P(h_c+u)-N\bigr).
\]
Here \(P\ge0\). In the nonnegative term replace \(\delta_c\) by
\(\delta_*\) and sum. The upper bound on \(\delta_c\) in
\eqref{eq:cov-shares} gives \(nB\delta_*\ge(1-\rho)W\). Hence
\[
 \sum a_c\ \ge\ B\delta_* Pnm-NW
       \ \ge\ W\bigl(P(1-\rho)m-N\bigr),\qquad 0<W\le S .
\]
As \(m\ge(1-\rho)u-X-v\),
\((1-\rho)m\ge(1-2\rho)u-X-v\). Subtract \(N\) and use
\(0\le P\le T\) on the terms with a minus sign. The result is
\[
 P(1-\rho)m-N\ \ge\ -T(X+v)-u(6D+10+T(z+2\rho))
              -A\log(\max(A,1))-4-R=:J .
\]
In fact the estimate before this last enlargement has \(6D+8\) in
place of \(6D+10\), no \(Tz\), and \(2\) in place of \(4\).
By construction \(J\le0\). If \(P(1-\rho)m-N\ge0\) then
\(\sum a_c/S\ge0\). Otherwise multiplication by \(W/S\le1\)
increases this negative quantity, so \(\sum a_c/S\ge J\) also in
that case. Substituting \eqref{eq:cov-budgets} and using \(4\le4E\)
proves \eqref{eq:cov-u}.
\footnote{\href{https://github.com/Deicyde/riemann-hypothesis/blob/40e56bab/src/Formalism/MeanHeightDescent/AverageDegreeBounds/SumComponents.lean}{Component sum};
\href{https://github.com/Deicyde/riemann-hypothesis/blob/40e56bab/src/Development/AbelianVarietyTheory/Proofs/DivisorProofs/UniformLowerBoundOfSums.lean}{specialization and common threshold}.}

At a given arithmetic point and a nonterminal lift index \(i\), the cover
supplies the finite set \(I\). For a quotient weight \(u_1\) the source
takes \(h_c\) and \(\delta_c\) at \(u_1\), and \(a_c,d_c\) at the
lower weight \(u_-\). In its notation, with \(P_0\) the size parameter,
\[
 T=\left\lceil P_0^2D^2 k/m_0\right\rceil,\qquad
 B=\binom{T-1}{D},\qquad z=7(D^2/T+D\rho),
\]
where \(D\) is the flag dimension, \(k\) is the weight parameter,
and \(m_0\) is the number of lifts. Here \(D,T\) are integers and
\(T\ge3D+5\) gives \(B>0\). One first chooses \(u_1\) among the
quotient weights guaranteed by the size condition. The degree and
rank-share estimates give \eqref{eq:cov-shares} for this same weight.
The component height estimate gives \eqref{eq:cov-shift} and the first
three budgets; the arithmetic Hilbert--Samuel estimate gives
\eqref{eq:cov-hs} and the fourth. The comparison for the sum of the
component degrees gives the last inequality of \eqref{eq:cov-budgets},
with \(\mu\) the cover slope at the weight indexed by \((u_-,0)\).
The seven uniform inputs each hold beyond a threshold. Their maximum
is a common threshold. Thus the constant \(C_U=2(3C_1+C_2+4)\)
works for every such cover on the same family.

\section{From the cover to the covolume of the theta lattice}
\label{sec:cov-lattice}

Fix a point, a lift index \(i\), and a cover at the required depth.
Write \(M\) for its main term, \(E\) for its error unit and \(b\) for
the scale in the norm bounds. Put \(K=P_0\kappa_i\), where \(\kappa_i\)
is the weight size. The size bounds include \(1\le b\le E,\ b\le K\)
and \(|M|\le K\). Write \(L_-\) for the cover lattice at the negative
shift and \(L_d\) for the integral lattice at the target weight
intersected with the prescribed type. Denote their squared-metric
slopes by \(\mu_-\) and \(\mu_d\); this normalization gives degree
equal to rank times slope, using squared norms at complex places.
Let \(\mu_0\) be the core slope at the original point.

The comparison for the flat line together with \eqref{eq:cov-u}
gives \(\mu_-\ge-M-(C_F+C_U)E\), with \(C_F\) the constant in the
flat-line comparison. The maps from \(L_-\) to \(L_d\) are integral,
have zero common kernel, and satisfy a sum-of-squares norm bound
with multiplier \(e^{2b}\). If there are \(j\) maps, the positive
shift bounds \(j\,\mathrm{rk}(L_d)\) above, and its rank is at most
\((1+b/K)\mathrm{rk}(L_-)\). The negative shift has positive rank.
The target is spanned over the field by integral sections of norm
at most \(e^K\) at each embedding. These are the hypotheses of the
integral lattice descent estimate in the source. With \(C=C_F+C_U\)
and \(t=-M-CE-2b\), that estimate gives
\[
                   \mu_d\ge t-(b/K)(|t|+2K).
\]
The comparison at the target ring gives \(\mu_0\ge\mu_d-2b\).
To see the exact loss, set \(\xi=b/K\). Then
\(\xi\le1,\ \xi K=b\), and
\[
 \xi |t|\le \xi|M|+\xi(CE+2b)\le b+CE+2b .
\]
Consequently \(\mu_0\ge-M-2CE-9b\ge-M-(2C+9)E\).
At the terminal lift there is a separate defect bound with constant
\(C_D\). Taking \(C_W=\max(C_D,2(C_F+C_U)+9)\) yields at every index
\begin{equation}                 \mu_0\ge-M-C_WE .\label{eq:cov-low}\end{equation}
\footnote{\href{https://github.com/Deicyde/riemann-hypothesis/blob/40e56bab/src/Development/AbelianVarietyTheory/Proofs/DivisorProofs/LowerBoundOfAllEstimates.lean}{Cover lattices and the numerical loss};
\href{https://github.com/Deicyde/riemann-hypothesis/blob/40e56bab/src/Development/AutomorphicArithmetic/ThetaCorrespondence/ThetaPairLifts/ThirdPlaceSeries/LowOfCoverHypotheses.lean}{last-index alternative}.}

\subsection{Removing the discarded subspace}
Let \(F\) be totally complex, and let \(L\) be a nonzero finitely
generated lattice in a hermitian \(F\)-space. Suppose \(e^2=e\), and at
each complex embedding squared norms add for \(ex\) and \((1-e)y\).
In the Gram realization write the inner product as conjugate-linear in
its first entry. The norm identity first gives
\(\Re\langle ex,(1-e)y\rangle=0\). At the chosen embedding pick
\(a\in F\) with nonreal image, which exists because the place is complex.
Applying the same identity with \(ax\) in place of \(x\) gives
\(\Re(\overline{a}\langle ex,(1-e)y\rangle)=0\). These two equations
force the inner product to vanish. Thus the image and kernel are
orthogonal and \(e\) is self-adjoint.

Let \(L_g=L\cap\mathrm{im}(e)\) and \(L_b=L\cap\ker(e)\), of ranks
\(d_g,d_b\). Suppose \(g\ge1\) is an integer such that \(geL\subset L\).
The map \(g(1-e)\) carries \(L\) into \(L_b\), its image \(J\)
contains \(gL_b\), and its kernel on \(L\) is \(L_g\). Hence
\(\operatorname{rk}L=d_g+d_b\) and \(J\) has rank \(d_b\).
The hermitian quotient by the saturated kernel identifies isometrically
at each embedding with the orthogonal projection onto \(\ker e\),
before multiplication by \(g\). Equivalently the operator from the
quotient to the image has norm at most \(g\). Taking determinants and
using \(J\subset L_b\) of full rank gives the cut inequality
\begin{equation}
 (d_g+d_b)\,\mu(L)\le d_g\,\mu(L_g)+d_b\{\mu(L_b)+2\log g\}.
 \label{eq:cov-cut}
\end{equation}
Indeed the squared archimedean determinant incurs at most \(g^{2d_b}\)
per place; finite index in \(L_b\) can only decrease the degree of
the image. Additivity of determinant degree in the saturated sequence
for \(L_g\subset L\) gives \eqref{eq:cov-cut}. Products of rank and
slope here mean degree, also on a rank-zero summand.

Fix \(0\le\delta\le1/12\) and suppose \(d_b\le3\delta(d_g+d_b)\).
Let \(Q\) bound the maximal slope of an ambient lattice containing \(L\),
as well as \(2\log g\) and \(-\mu(L)\). Then \(Q\ge0\),
\(d_g>0\), and \(d_b\le4\delta d_g\) since \(1-3\delta\ge3/4\).
In \eqref{eq:cov-cut} the bad degree is at most \(d_bQ\) by the
definition of maximal slope. Moving its terms to the left gives
\[
 d_g(\mu(L)-\mu(L_g))\le d_b(Q+2\log g-\mu(L))
                 \le 3d_b Q\le12\delta d_gQ .
\]
This proves \(\mu(L_g)\ge\mu(L)-12\delta Q\), including the case \(d_b=0\).
\footnote{\href{https://github.com/Deicyde/riemann-hypothesis/blob/40e56bab/src/Formalism/MeanHeightDescent/AverageDegreeBounds/GoodBadCut.lean}{Integral cut and height of the image};
\href{https://github.com/Deicyde/riemann-hypothesis/blob/40e56bab/src/Development/AbelianVarietyTheory/Proofs/DivisorProofs/MajAtAddRight.lean}{orthogonality and the factor twelve}.}

\subsection{The uniform conclusion}
In the application \(L\) is the lattice of the prescribed type and
\(L_g\) its retained part. Put \(\delta=\exp(-x^{\varepsilon_0/7})\),
where \(\varepsilon_0>0\) is fixed on the family. The truncation
provides the idempotent above, \(g\) the order of its finite group,
and \(2\log g\le b\). The discarded rank is at most \(3\delta\,\mathrm{rk}L\).
Suppose the maximal slope is bounded by \(x^A b\), \(|M|\le x^B b\),
and \(\mu(L)\ge\mu_0-C_KE\) for \(C_K\ge0\). By \eqref{eq:cov-low},
\(-\mu(L)\le M+(C_K+C_W)E\). Then
\[
 Q=(x^A+x^B+C_K+C_W+1)E
\]
dominates each of the three bounds required above. The condition
\(12\delta(x^A+x^B+C_K+C_W+1)\le1\) also ensures \(\delta\le1/12\).
It holds eventually because \(x^j\exp(-x^{\varepsilon_0/7})\to0\)
for each real \(j\). Therefore \(\mu(L_g)\ge\mu(L)-E\).
The preceding recursion gives the normalized log covolume at \(i\)
at most \(-\mu(L_g)\). Combining the three inequalities yields
\[
  \frac{\log\operatorname{covol}_i}{D_{\rm norm}}
      \le M+(C_K+C_W+1)E .
\]
Here the left side denotes precisely the lattice log covolume and
normalization in the recursion at that index. All inequalities hold
simultaneously above the maximum of the loss threshold and the
thresholds in the eight uniform inputs. For each larger constant the
conclusion persists because \(E\ge1\).
\footnote{\href{https://github.com/Deicyde/riemann-hypothesis/blob/40e56bab/src/Development/AbelianVarietyTheory/Proofs/DivisorProofs/CovolumeFromGramAndBounds.lean}{The pointwise bound};
\href{https://github.com/Deicyde/riemann-hypothesis/blob/40e56bab/src/Development/AutomorphicArithmetic/ThetaCorrespondence/ThetaPairLifts/ThirdPlaceSeries/CovolXOfSteps.lean}{the uniform covolume conclusion}.}

The cover degree, its rank share, its component height and
Hilbert--Samuel bounds, and the recursive covolume comparison were
inputs in this derivation. Their formal constructions enter through
\href{https://github.com/Deicyde/riemann-hypothesis/blob/40e56bab/src/Development/AutomorphicArithmetic/ThetaCorrespondence/ThetaPairLifts/ThirdPlaceSeries/TopCovolG.lean}{the geometric assembly theorem}.
